\begin{frame}[fragile]{Step 1: Make a Module (sum.js)}

\begin{itemize}
	\item When we need a function that adds to argument values, we can implement the function in a module (sum.js). 
	\item We need to export the `sum' function using the `module.exports'.
\end{itemize}

\begin{Code}{}%
\begin{lstlisting}[firstnumber=1, label=glabels, xleftmargin=0pt, basicstyle=\scriptsize]% 

function sum(a, b) {
  return a + b;
}
module.exports = sum;
\end{lstlisting}%   
\end{Code}%

\end{frame}

\begin{frame}[fragile]{Step 2: Make a Unit Test (sum.test.js)}

\begin{itemize}
	\item In this example, we have the module and test in the same directory, so we use line 1 to access the function to test.
	\item We use Jest's test() function that has (a) the comment and (b) lambda expression.
	\item The body of lambda expression has two parts (a) expect() that has the function call and (b) toBe() that has the expected value.
\end{itemize}

\begin{Code}{}%
\begin{lstlisting}[firstnumber=1, label=glabels, xleftmargin=0pt, basicstyle=\tiny]% 

const sum = require('./sum');

test('adds 1 + 2 to equal 3', () => {
  expect(sum(1, 2)).toBe(3);
});
\end{lstlisting}%   
\end{Code}%

\end{frame}

\begin{frame}[fragile]{Step 3: Set Up Jest}

\begin{itemize}
	\item We need to add ``test'': ``jest'' in the ``script'' section of the package.json. 
	\item We instruct the npm command that when we invoke ``test'', the command ``jest'' is called. 
	\item We can add additional arguments to the jest command if necessary. 
\end{itemize}

\begin{Code}{}%
\begin{lstlisting}[firstnumber=1, label=glabels, xleftmargin=0pt, basicstyle=\scriptsize]% 

"scripts": {
  "test": "jest"
},
\end{lstlisting}%   
\end{Code}%

\end{frame}

\begin{frame}[fragile]{Run Jest}

\begin{itemize}
	\item We can use `npm test' to run Jest. 
	\item We also can execute jest command line with `jest'.
	\item Jest automatically calls the Jest test files to show the results. 
\end{itemize}

\begin{center}
\includegraphics[width=0.5\textwidth]{res.png}
\end{center} 
\end{frame}

\begin{frame}[fragile]{Run Single Jest Test}

\begin{itemize}
	\item When we run jest without any arguments, jest finds all the test files (such as the file that ends with `.test.js').
	\item When we want to run single test, we can use the `-i' option to specify the test that we want to run. 
\end{itemize}

\begin{Code}{}%
\begin{lstlisting}[firstnumber=1, label=glabels, xleftmargin=0pt, basicstyle=\scriptsize]% 

"scripts": {
  "test": "jest -i sum.test.js"
},
\end{lstlisting}%   
\end{Code}%

\begin{Code}{}%
\begin{lstlisting}[firstnumber=1, label=glabels, xleftmargin=0pt, basicstyle=\scriptsize]% 

jest -i sum.test.js
\end{lstlisting}%   
\end{Code}%

\end{frame}

\begin{frame}[fragile]{When We Find Bugs}

\begin{itemize}
	\item We may find a bug that can pass tests.
	\item In this case, we should (a) fix the bug and (b) add unit tests to catch the missed bug. 
\end{itemize}
\end{frame}